% Generated draft: requires mathematical and source review.
\documentclass[11pt]{article}
\usepackage[margin=25mm]{geometry}
\usepackage{fontspec}
\setmainfont{Latin Modern Roman}
\usepackage{amsmath,amssymb,mathtools,unicode-math}
\setmathfont{Latin Modern Math}
\usepackage{enumitem,needspace,xurl,hyperref}
\setlength{\parindent}{0pt}
\setlength{\parskip}{5pt}
\setlength{\emergencystretch}{3em}
\newcommand{\N}{\mathbb{N}}
\newcommand{\Z}{\mathbb{Z}}
\newcommand{\Q}{\mathbb{Q}}
\newcommand{\R}{\mathbb{R}}
\newcommand{\C}{\mathbb{C}}
\newcommand{\F}{\mathbb{F}}
\newcommand{\A}{\mathbb{A}}
\newcommand{\PP}{\mathbb{P}}
\newcommand{\E}{\mathbb{E}}
\newcommand{\eps}{\varepsilon}
\newcommand{\dd}{\,\mathrm{d}}
\begin{document}
\section*{UnsolvedMath ID: 30002175}
\textbf{Source problem code:} catalogue entry
\subsection{Definitions and mathematical statement}
Under the conventional interpretation, let \(S_{g,n}\) be a connected orientable surface of genus \(g\in\Z_{\geq0}\) with \(n\in\Z_{\geq0}\) punctures and no boundary. This interpretation is provisional: the supplied record does not specify the surface or mapping-class conventions, including whether punctures must be fixed individually. Fixing those conventions, let \(\mathcal P_{g,n}\) denote the set of pseudo-Anosov mapping classes on \(S_{g,n}\). A pseudo-Anosov representative preserves a pair of transverse measured singular foliations, scaling their transverse measures by factors \(\lambda\) and \(\lambda^{-1}\), where \(\lambda>1\). Its dilatation is \(\lambda\). Whenever \(\mathcal P_{g,n}\) is nonempty and a least dilatation exists, write
\[
\delta(S_{g,n})=\min\{\lambda(f):f\in\mathcal P_{g,n}\}.
\]
The source does not specify the admissible pairs \((g,n)\) or a convention for cases with no pseudo-Anosov mapping classes; no value of \(\delta\) is assigned here in those cases.

The supplied problem asks to determine the dependence of \(\delta(S_{g,n})\) on genus and puncture number, and to describe mapping classes with small dilatation, including the structure of minimizers mentioned in the supplied background. It gives no quantitative meaning for ``small,'' no required form of the dependence, and no regime in which \(g\) and \(n\) vary.

It also asks whether normalized dilatations have an accumulation point smaller than
\[
\ell_0=\left(\frac{3+\sqrt5}{2}\right)^2.
\]
The normalization and the collection over which accumulation is taken are not supplied. Conditionally, if \(D\subset\R\) is the intended set of normalized dilatation values and ``accumulation point'' has its usual meaning for a subset of \(\R\), this question is
\[
\exists L\in\R\quad L<\ell_0\quad\text{and}\quad
\forall\varepsilon>0\ \exists d\in D:\ 0<|d-L|<\varepsilon.
\]
This is a conditional formulation only: the record does not determine \(D\), whether it includes all pseudo-Anosov classes or only minimum dilatations, or whether limits along distinct surfaces or mapping classes are intended instead. These qualifications require resolution before the accumulation question has a uniquely specified mathematical target.
\subsection{Short English statement}
Determine how the least pseudo-Anosov dilatation depends on genus and punctures, and describe classes with small dilatation and the minimizers. Decide whether normalized dilatations accumulate below \(\ell_0=((3+\sqrt5)/2)^2\). The supplied record leaves the normalization, accumulation family, surface conventions, and meaning of ``small'' unspecified, so this remains a qualified statement-only draft.
\subsection{Sources}
\begin{itemize}
\item Supplied local UnsolvedMath registry record, permanent ID \(30002175\), title \emph{Minimal Pseudo-Anosov Dilatations}, source problem code \texttt{catalogue entry}, in the Oberwolfach Reports Open Problems collection. No individual proposer is identified; the proposed year is recorded as 2012.
\item Supplied report attribution: \emph{Low-Dimensional Topology and Number Theory} (2013), Oberwolfach Reports, pp.\ 2541--2596, DOI \texttt{10.4171/owr/2012/42}. \url{https://doi.org/10.4171/owr/2012/42}. The background describes these reports as records of mathematical workshops and their open-problem discussions. The literature-source label is ``OWR report 42.''
\item Supplied literature source labeled ``Hironaka'': \url{https://arxiv.org/abs/1403.2987}. No title or specific supporting result is supplied.
\item Supplied literature source labeled ``Valdivia'': \url{https://nyjm.albany.edu/j/2012/18-32v.pdf}. No title or specific supporting result is supplied.
\item Literature assessment dated exactly \texttt{2026-08-21}: the supplied curation labels the problem ``partially solved'' and reports sharp asymptotic bounds, explicit small-dilatation families, and further results for restricted classes, while describing the full genus--puncture dependence, the smallest accumulation point, and the structure of all minimizers as open. This is supplied dated curation, not an independently verified assessment. The background explicitly calls for independent verification of time-sensitive status and cited claims.
\item The registry labels the statement \texttt{corrected\_verified} and records a manual semantic, typographic, and LaTeX review of its cleanup. Those supplied labels do not establish mathematical completeness or independent source verification here.
\end{itemize}
\end{document}
